% R038：复用例1数据 (1,2),(2,3),(3,5),(4,6)，中心为 (2.5,4)。
\begin{center}
\begin{tikzpicture}[x=1cm,y=1cm,>=Stealth,every node/.style={font=\small}]
\path[use as bounding box] (0,0) rectangle (12.4,6.35);
\node[font=\large\bfseries,text=CProof] at (6.2,6.02) {中心化：样本中心移到新原点，样本之间的相对位置不变};

\draw[rounded corners=5pt,fill=purple!2,draw=CProof] (.18,1.25) rectangle (5.82,5.42);
\node[font=\bfseries,text=CProof] at (3.00,5.04) {原坐标：例1的四个样本点};
\draw[->,black!65] (.82,1.70)--(5.32,1.70) node[right] {$x$};
\draw[->,black!65] (.82,1.70)--(.82,4.95) node[above] {$y$};
\foreach \x/\y in {1.40/2.70,2.40/3.20,3.40/4.20,4.40/4.70}{\fill[blue!80!black] (\x,\y) circle (2.8pt);}
\fill[orange!90!black] (2.90,3.70) circle (4.0pt);
\node[anchor=south,text=orange!85!black,fill=purple!2,inner sep=1pt] at (2.90,3.76) {$(\bar x,\bar y)=(2.5,4)$};

\draw[->,CProof,very thick] (6.03,3.63)--(6.43,3.63) node[midway,above=3pt,font=\scriptsize] {平移};
\node[align=center,text=CProof,font=\scriptsize] at (6.23,2.98) {$(x_i,y_i)\mapsto$\\$(x_i-\bar x,y_i-\bar y)$};

\draw[rounded corners=5pt,fill=purple!2,draw=CProof] (6.58,1.25) rectangle (12.22,5.42);
\node[font=\bfseries,text=CProof] at (9.40,5.04) {中心化后：新坐标轴经过 $O$};
\coordinate (O) at (9.40,3.70);
\draw[->,black!65] (7.18,3.70)--(11.85,3.70) node[right] {$x'$};
\draw[->,black!65] (9.40,1.72)--(9.40,4.95) node[above] {$y'$};
\foreach \x/\y in {7.90/2.70,8.90/3.20,9.90/4.20,10.90/4.70}{\fill[blue!80!black] (\x,\y) circle (2.8pt);}
\draw[orange!90!black,line width=1.1pt] (O) circle (4.0pt);
\node[below left,text=orange!85!black] at (O) {$O=(0,0)$};

\node[align=center,text=CProof] at (6.2,.66) {中心化后 $\sum x_i'=0,\ \sum y_i'=0$；橙色的样本中心准确移到 $O$，\textbf{它不是原始样本点。}};
\end{tikzpicture}
\end{center}
